\section{All pre-registered verdicts}\label{app:verdicts}

Table~\ref{tab:verdicts} lists the \nHypotheses{} hypotheses of the five pre-registrations, each with its pre-registered
verdict. A superiority verdict counts only if its rule passes at most 5\% of label-shuffled copies; this removed
CP4. For the hypotheses that hold, $p$ is that of the rule's primary contrast (for a rule that needs $k$ passing
budgets, the $k$-th smallest); all \nHolm{} survive Holm's correction~\cite{holm1979} over all \nHypotheses. The full
rules, and every change after a freeze with its date and reason, are in the released \texttt{prereg/} folder.

\begin{table*}
  \caption{The pre-registered hypotheses and their verdicts.}\label{tab:verdicts}
  \small\verdictTable
\end{table*}

\section{Parser sensitivity}\label{app:parser}

Runs set aside when the parser changed, compared with the final run on the same \texttt{base} tasks
(Table~\ref{tab:parser}). Quality is the macro pass rate over tracks; $\Delta$ is final minus old.

\begin{table*}
  \caption{Old-parser runs against the final parser on the same tasks.}\label{tab:parser}
  \small\parserTable
\end{table*}

\section{Calibration}\label{app:calibration}

For each superiority rule we built 50 copies of the data in which, per task, the labels of the compared systems are
exchanged at random, and reran the whole selection and test. A rule that passes more than 5\% of the copies is
treated as not shown. Two problems were found this way and are dated in the pre-registration logs. In stage 2, the
switch-versus-retry cascade (HR3) passed 48\% of copies: each policy's detector threshold was chosen on the test
tasks it was scored on, and the two policies were compared at their own best thresholds, so they escalated
different shares of tasks to the large model and that difference, not the retry, decided the test. With thresholds
chosen on fit data and both policies at the retry's threshold, the rule passes \switchFalse{} of copies and still
holds. In stage 4, the rule ``the router beats the best single model at any of 16 budgets'' passed \cpFourFalse{} of
copies, so CP4 is not shown although its best budget gains \routerVsSingle\pp. CP3 is a non-inferiority claim, which
equal systems satisfy, so the shuffle does not test it. CP5 could not be evaluated: the tiny stack is cheap because
its episodes end early, and no router fits its cost.
